\documentclass{article}
\usepackage[utf8]{inputenc}
\usepackage{booktabs}
\usepackage{caption}
\usepackage{threeparttable}
\usepackage{amsmath}
\usepackage{pdflscape}

\begin{document}

\begin{landscape}
\begin{table}[htbp]
    \centering
    \begin{threeparttable}
    \caption{Comparative Balance between Urban and Rural Subgroups}
    \label{tab:balance_subgroups}
    \begin{tabular}{l ccc ccc}
        \toprule
        & \multicolumn{3}{c}{\textbf{Urban Subsample}} & \multicolumn{3}{c}{\textbf{Rural Subsample}} \\
        \cmidrule(lr){2-4} \cmidrule(lr){5-7}
        & {Control} & {Treated} & {$p$-val} & {Control} & {Treated} & {$p$-val} \\
        Variable & {(N=650)} & {(N=642)} & {(diff)} & {(N=550)} & {(N=543)} & {(diff)} \\
        \midrule
        Age (years)      & 34.5  & 34.2  & 0.84 & 38.2  & 38.5  & 0.72 \\
                         & (8.2) & (8.1) &      & (9.1) & (9.0) &      \\
        Female (\%)      & 0.52  & 0.51  & 0.91 & 0.48  & 0.49  & 0.81 \\
                         & (0.5) & (0.5) &      & (0.5) & (0.5) &      \\
        Income (\$k)     & 52.4  & 53.1  & 0.65 & 25.1  & 24.8  & 0.78 \\
                         & (12.4)& (12.8)&      & (8.5) & (8.2) &      \\
        Education (yrs)  & 14.2  & 14.1  & 0.88 & 8.4   & 8.5   & 0.92 \\
                         & (2.4) & (2.3) &      & (3.1) & (3.0) &      \\
        \midrule
        \textbf{Joint F-test (p)} & & & 0.95 & & & 0.88 \\
        \bottomrule
    \end{tabular}
    \begin{tablenotes}
        \small
        \item \textit{Notes:} This comparative balance table ensures randomization was successful within key population subgroups (Urban vs Rural). Standard deviations in parentheses.
    \end{tablenotes}
    \end{threeparttable}
\end{table}
\end{landscape}

\end{document}
